\chapter{Spleen}

\label{ch:ch48_spleen}

\section{What does the spleen do?}
The spleen sits under your left ribcage and acts as a blood filter. It removes old or damaged red blood cells, stores platelets, and helps fight certain infections. Because it is rich in blood vessels, injuries to the spleen can cause significant bleeding, which is why doctors take abdominal pain on the left side seriously, especially after trauma.

\section{When is splenic imaging ordered?}
Doctors may order imaging of the spleen for several reasons:
\begin{itemize}
    \item \textbf{Trauma:} After a car accident, fall, or blow to the abdomen, imaging checks whether the spleen is injured or bleeding.
    \item \textbf{Unexplained left upper quadrant pain:} This may suggest infarction, infection, or a focal lesion.
    \item \textbf{Enlarged spleen found on examination:} Ultrasound or CT measures the spleen and looks for a cause.
    \item \textbf{Known cancer or lymphoma:} Imaging stages disease and assesses splenic involvement.
    \item \textbf{Fever of unknown origin:} A splenic abscess or granulomatous disease may be identified.
\end{itemize}

\section{What happens during the scan?}
\begin{itemize}
    \item \textbf{Ultrasound:} A handheld probe is moved over the left side of the abdomen with gel. No radiation is used. The scan takes 10--15 minutes. An IV line may be placed if Doppler (blood flow) assessment is needed.
    \item \textbf{CT scan:} You lie on a table that slides through a ring-shaped scanner. A contrast injection may be used to improve detail. The scan takes about 5 minutes. You may feel a warm flush from the contrast---this is normal.
    \item \textbf{MRI:} A longer scan (20--40 minutes) in a narrow tube. It uses a strong magnet rather than radiation and is helpful for characterizing focal splenic lesions.
\end{itemize}

\section{Normal splenic appearance}
The normal spleen is homogeneous (uniform in texture) on ultrasound and CT. It is crescent-shaped, smooth-contoured, and typically measures up to 12~cm in craniocaudal length on ultrasound. On CT, the splenic parenchyma enhances uniformly after contrast administration. The splenic artery and vein are visible at the hilum.

\section{What might the report say?}
\begin{itemize}
    \item \textbf{Splenomegaly:} The spleen is larger than normal. Common causes include infection (mononucleosis, malaria), liver disease with portal hypertension, blood disorders (sickle cell disease, thalassaemia), and malignancy (lymphoma, leukaemia).
    \item \textbf{Focal lesions:} These are bumps or areas within the spleen. Most are benign---cysts, haemangiomas, and hamartomas are common incidental findings. Lymphoma is the most common malignant cause of splenic lesions.
    \item \textbf{Laceration or haematoma:} A tear in the spleen with or without surrounding blood, typically after trauma.
    \item \textbf{Infarction:} A wedge-shaped area of non-enhancement, often due to embolism or a hypercoagulable state.
    \item \textbf{Accessory spleen:} A small nodule of normal splenic tissue near the main spleen---a normal variant, not a disease.
\end{itemize}

\section{Splenic size: what is normal?}
Table~\ref{tab:spleen_size} summarises reference values.

\begin{table}[h]
\centering
\caption{Splenic size reference values}
\label{tab:spleen_size}
\begin{tabular}{lll}
\hline
\textbf{Measurement} & \textbf{Normal} & \textbf{Splenomegaly} \\
\hline
Anteroposterior diameter (US) & $<$5~cm & $\geq$5~cm \\
Craniocaudal length (US) & $<$12~cm & $\geq$12~cm \\
Volume (CT/MRI) & $<$200~mL & $>$200~mL \\
\hline
\end{tabular}
\end{table}

\section{Splenic trauma and the AAST grading system}
The American Association for the Surgery of Trauma (AAST) grading system guides management of splenic injuries:

\begin{table}[h]
\centering
\caption{AAST splenic injury grading}
\label{tab:aast_spleen}
\begin{tabular}{clll}
\hline
\textbf{Grade} & \textbf{Description} & \textbf{Management} & \textbf{Mortality} \\
\hline
I & Subcapsular haematoma $<$10\% surface area; laceration $<$1~cm depth & Observation & $<$1\% \\
II & Subcapsular haematoma 10--50\%; laceration 1--3~cm depth & Observation & 1--2\% \\
III & Subcapsular haematoma $>$50\% or expanding; laceration $>$3~cm depth & Observation or angio & 3--5\% \\
IV & Laceration involving segmental vessels or hilum & Angioembolisation $\pm$ surgery & 10--20\% \\
V & Shattered spleen or hilar vascular injury & Emergency surgery & $>$50\% \\
\hline
\end{tabular}
\end{table}

\section{When to recommend angiography, surgery, or observation}
\begin{itemize}
    \item \textbf{Observation:} Appropriate for Grades I--III in haemodynamically stable patients. Serial abdominal examinations and repeat imaging at 24--48 hours assess for interval haemorrhage.
    \item \textbf{Angiography with embolisation:} Indicated for Grade III--IV injuries with contrast blush (active extravasation) on CT, or for patients with ongoing transfusion requirements despite initial stability. Proximal embolisation preserves splenic function while controlling bleeding.
    \item \textbf{Surgery (splenectomy):} Reserved for haemodynamically unstable patients failing resuscitation, Grade V injuries, or failed angiographic embolisation. Partial splenectomy may be considered in selected cases.
\end{itemize}

\begin{pearlbox}[Scalloping and lymphoma]
Multiple small, rounded low-density splenic lesions (``Swiss cheese'' pattern) on CT or multiple T2-dark nodules on MRI should raise suspicion for lymphoma. Isolated splenic lymphoma is uncommon; staging imaging of the chest, abdomen, pelvis, and bone marrow is recommended.
\end{pearlbox}

\section{Focal splenic lesions: a diagnostic approach}
Table~\ref{tab:splenic_lesions} guides lesion characterisation.

\begin{table}[h]
\centering
\caption{Common focal splenic lesions}
\label{tab:splenic_lesions}
\begin{tabular}{llll}
\hline
\textbf{Lesion} & \textbf{US} & \textbf{CT} & \textbf{Key feature} \\
\hline
Abscess & Complex, internal debris & Rim enhancement, gas & Fever, clinical correlation \\
Haemangioma & Hyperechoic & Peripheral nodular enhancement & Most common solid lesion \\
Hamartoma & Iso- to hyperechoic & Hypodense, heterogeneous enhancement & May mimic lymphoma \\
Infarction & Hypoechoic, wedge-shaped & Wedge-shaped non-enhancement & Pain, risk factors \\
Lymphoma & Hypoechoic, multiple & Hypodense, mild enhancement & Stage with PET/CT \\
Simple cyst & Anechoic, posterior enhancement & Water density, no enhancement & Benign, no treatment \\
\hline
\end{tabular}
\end{table}

\section{Accessory spleen}
An accessory spleen (splenule) is found in approximately 10--30\% of people. It is a small nodule of normal splenic tissue, typically 1--3~cm, located near the splenic hilum. It enhances identically to the main spleen on all phases of contrast-enhanced imaging. Its importance lies in:
\begin{itemize}
    \item Being mistaken for a lymph node or pathological mass.
    \item Persistence after splenectomy for haematological disease (causing treatment failure in conditions such as immune thrombocytopenia).
\end{itemize}

\section{Splenic infarction}
Splenic infarction presents with acute left upper quadrant pain. Causes include:
\begin{itemize}
    \item Cardiac embolism (atrial fibrillation, endocarditis)
    \item Haematological disorders (sickle cell disease, myeloproliferative neoplasms)
    \item Vasculitis and hypercoagulable states
    \item Iatrogenic (post-procedure, splenic artery ligation)
\end{itemize}
On imaging, infarction appears as a wedge-shaped, non-enhancing area with its base at the splenic capsule. Over time, the area may cavitate or calcify.

\section{Reporting checklist}
A structured splenic report should include:
\begin{enumerate}
    \item Spleen size (measurements)
    \item Parenchymal echogenicity/density pattern
    \item Focal lesions: number, size, location, enhancement pattern
    \item Trauma findings: laceration depth, subcapsular haematoma volume, active extravasation
    \item Perisplenic fluid
    \item Vascular anatomy (splenic artery/vein patency)
    \item Accessory spleen
    \item Comparison with prior imaging
\end{enumerate}

\clearpage
